% ECONOMICS_PROBLEM: {"id": "J14-452", "title": "Healthy aging and demographic burdens", "jel_code": "J14", "source_id": "OP-0459"}
% !TeX program = xelatex
\documentclass[11pt,a4paper]{article}
\usepackage[margin=23mm]{geometry}
\usepackage{fontspec}
\setmainfont{Latin Modern Roman}
\usepackage{amsmath,amssymb,mathtools}
\usepackage{xcolor,enumitem,booktabs,longtable,array,xurl,hyperref}
\definecolor{muted}{HTML}{53636C}
\hypersetup{colorlinks=true,urlcolor=blue,linkcolor=blue}
\setlength{\parindent}{0pt}
\setlength{\parskip}{5pt}
\newcommand{\R}{\mathbb R}
\newcommand{\E}{\mathbb E}
\newcommand{\N}{\mathbb N}
\newcommand{\ind}{\mathbf 1}
\newcommand{\Var}{\operatorname{Var}}
\newcommand{\Cov}{\operatorname{Cov}}
\newcommand{\supp}{\operatorname{supp}}
\DeclareMathOperator*{\argmin}{arg\,min}
\DeclareMathOperator*{\argmax}{arg\,max}
\newcommand{\entrymeta}[3]{{\sffamily\small\color{muted}\textbf{\texttt{#1}}\quad|\quad #2\par #3\par}}
\newcommand{\Problemref}[1]{\texttt{#1}}
\newcommand{\JELref}[2]{\texttt{#2} (JEL \texttt{#1})}
\begin{document}
\section*{Healthy aging and demographic burdens}
\textbf{Problem ID:} \texttt{J14-452}\quad \textbf{JEL:} \texttt{J14}\par
% BEGIN ECONOMICS STATEMENT
\label{problem:OP-0459}
{\sffamily\small\color{muted}Original subject group: Education, families and demography\par}
\label{definitions:problem:30007741}
\textbf{Audit classification:} Background; proposed extension. Specified survival accounting and separated fiscal pension transfers from real welfare costs.
\subsection*{Definitions and precise research target}
\textit{Editorial formalization.} Consider residents aged fifty through eighty-five in one high-income country over a twenty-year fiscal projection. Let \(a=1\) denote a specified improvement in age-specific functional health relative to baseline \(a=0\), preserving a stated mortality path or explicitly modeling its change. For individual \(i\), define annual labor earnings \(Y_{it}(a)\), public health and care expenditure \(H_{it}(a)\), pension transfers \(P_{it}(a)\), and taxes \(T_{it}(a)\). For discount factor \(\beta\in(0,1)\), define net fiscal burden
\[B(a)=E\!\left[\sum_i\sum_{t=0}^{20}\beta^t\{H_{it}(a)+P_{it}(a)-T_{it}(a)\}\right].\]
Determine \(B(1)-B(0)\) under stated retirement rules, wage responses, caregiving arrangements, and demographic projections. Estimate health’s effects on work and care needs rather than infer behavior from chronological age alone. Separate healthier aging from longer survival, and report distributional household welfare alongside the fiscal comparison. Causal identification must address selection in health and employment; observational age profiles are insufficient.

Fix the baseline cohort aged 50–85 and a twenty-year projection date. Let \(s_{it}(a)\) be survival; costs, pensions and taxes after death are zero unless survivor benefits are separately allocated. Health intervention \(a\) specifies an age/cohort functional-health path rather than changing every unmeasured characteristic. \(B(a)\) uses actual government outlays and receipts at common prices, so pension transfers belong in a fiscal burden even though they are not real social costs. Labor earnings \(Y\) enter through taxes and behavior, not again as a separate budget receipt. If longevity changes, population weights and future births/immigrants follow a declared demographic accounting rule. Finite expected discounted flows and explicit retirement, wages and care equations define the contrast. It is fiscal, not a complete welfare valuation. All expectations use a stated baseline population and probability law. Identification means every admissible data-generating model with the same observed-data law yields the same displayed target; otherwise report the identified set. Exact equations, horizons and assumptions are editorial, rather than source-given canonical conjectures.
\subsection*{Variants and assumption changes}
\begin{enumerate}
\item \textbf{Editorial longevity variant.} Allow health to change survival and jointly estimate fiscal burden and quality-adjusted years. Longer life may raise pension outlays.
\item \textbf{Editorial retirement variant.} Cross health improvement and retirement-age reform; estimate \(B(1,1)-B(1,0)-B(0,1)+B(0,0)\) including partner labor and caregiving.
\end{enumerate}
\subsection*{References and source audit}
\textbf{1.} Age versus physiological function study; exact fiscal-offset target is editorial. Locator: IZA Discussion Paper 17427, October 2024; abstract. Full text or primary publication page inspected. URL: \url{https://docs.iza.org/dp17427.pdf}.
\begin{enumerate}\raggedright
\item\hypertarget{definitions:source:30007741:1}{} Rainer Kotschy; David E. Bloom; Andrew J. Scott. 2024. \emph{On the Limits of Chronological Age}. IZA Discussion Paper 17427, October 2024; abstract.
\url{https://docs.iza.org/dp17427.pdf}.
\end{enumerate}

\subsection*{Common conventions: Source mathematical conventions}

These conventions apply to each entry unless explicitly replaced.

\textbf{Models and identification.} A primitive model \(m\) specifies all probability laws, feasible choices, information, equilibrium and measurement rules used by its target. The admissible class \(\mathcal M\) is fixed independently of outcomes. If \(P_O(m)\) is its observable probability law and \(\tau(m)\) the target, its identified set at observable law \(P\) is
\[
 \mathcal I_\tau(P)=\{\tau(m):m\in\mathcal M,\ P_O(m)=P\}.
\]
Point identification means that this set is a singleton. An empty set signals incompatibility of data and assumptions. Coordinate bounds need not describe the joint set. Bounds are sharp only if every retained target is attainable or arbitrarily approachable by compatible models. Finite bounds require explicit restrictions on unobserved quantities; unrestricted confounding, hidden wealth, extrapolation or arbitrary equilibrium selection can yield uninformative sets.

\textbf{Causal interventions.} A potential outcome \(Y_i(p)\) is the outcome under a complete policy regime \(p\), including timing, financing, information and market spillovers. The target population and units are fixed before treatment. With interference, use the complete assignment vector or a justified exposure mapping; individual no-interference notation alone cannot identify an economy-wide intervention. A difference of mean potential outcomes does not require the cross-world joint distribution. Conditioning on post-treatment migration, survival, enrollment or employment changes the estimand and needs separate justification.

\textbf{Statistical statements.} An estimator, confidence set, forecast or test specifies a sampling law, model class and loss/coverage criterion. Uniform \((1-\alpha)\)-coverage means \(\inf_{m\in\mathcal M}\Pr_m\{\tau(m)\in C_n\}\geq1-\alpha\), or an explicitly asymptotic version. The whole parameter space trivially has coverage, so informativeness requires a stated width, risk or power objective. A forecasting improvement is relative to a named benchmark, information set and evaluation population.

\textbf{Equilibrium and optimization.} An equilibrium comprises feasible plans, optimality, beliefs where needed, budget constraints and all market/resource clearing. The primitives or a stated benchmark define each correspondence \(\mathcal E(p;m)\). Nonemptiness is assumed or proved, never inferred just from compact policy sets. A selected-equilibrium target uses a declared selection rule; a robust target quantifies over all permitted equilibria. Use suprema or infima unless attainment is established. An \(\varepsilon\)-optimal policy is feasible and within \(\varepsilon\) of the finite optimum value. Report an empty feasible set separately.

\textbf{Welfare and accounting.} Normative utility, interpersonal normalization, social weights, numeraire, discounting and the use of public receipts are primitives. Behavioral utility need not coincide with the welfare criterion. Household welfare includes ownership income and rents once; adding profits again to compensating variation generally double-counts them. A monetary equivalent is defined by an explicit transfer and price convention, with monotonicity and range conditions. Average effects, mechanism contributions, transfers and welfare losses are different targets. Infinite sums require convergence and dynamic feasible plans require terminal or no-Ponzi conditions.

\textbf{Source versions.} A working-paper date, revision date and journal publication year are distinguished. A DOI for a journal version is not automatically the DOI of an earlier manuscript. Locators refer to the stated version and printed pages where specified. A metadata-only check does not verify a claim about a full-text conclusion. Every entry's source subsection narrows the provenance claim accordingly.
% END ECONOMICS STATEMENT
\end{document}
